% Slides for the write-up writeup-2026-11-15T0054Z.tex.
% This is the template: run `python3 make_slides.py` in this folder to fill the %%...%% markers
% with the pictures and numbers it computes, and compile the .tex file it writes.
\documentclass[aspectratio=169,11pt,t]{beamer}
\usetheme{default}
\usefonttheme{professionalfonts}
\usepackage[T1]{fontenc}
\usepackage{tgheros}
\renewcommand{\familydefault}{\sfdefault}
\usepackage{booktabs}
\setbeamertemplate{navigation symbols}{}
\setlength{\parskip}{0.45em}

% the style of the team's earlier slides: warm paper, dark plum ink, orange accent
\definecolor{paper}{HTML}{FBF8F4}
\definecolor{ink}{HTML}{2B0B1E}
\definecolor{accent}{HTML}{F26322}
\definecolor{muted}{HTML}{79646E}
\definecolor{rulecol}{HTML}{E1DBD9}
\newcommand{\eventline}{Caltech $\times$ SAIR Mathathon 2026}
\setbeamercolor{background canvas}{bg=paper}
\setbeamercolor{normal text}{fg=ink}
\setbeamercolor{frametitle}{fg=ink}
\setbeamercolor{framesubtitle}{fg=ink}
\setbeamercolor{title}{fg=ink}
\setbeamercolor{structure}{fg=accent}
\setbeamercolor{alerted text}{fg=accent!85!black}
\setbeamercolor{block title}{fg=ink, bg=accent!16}
\setbeamercolor{block body}{fg=ink, bg=white}
\setbeamercolor{block title example}{fg=ink, bg=rulecol}
\setbeamertemplate{blocks}[rounded][shadow=false]
\setbeamerfont{frametitle}{size=\large, series=\bfseries}
\setbeamerfont{framesubtitle}{size=\footnotesize, series=\bfseries}
\setbeamerfont{title}{size=\LARGE, series=\bfseries}
\setbeamertemplate{itemize items}{\textbullet}
\makeatletter
\setbeamertemplate{frametitle}{%
  \vspace*{0.6em}%
  {\usebeamerfont{frametitle}\usebeamercolor[fg]{frametitle}\insertframetitle\par}%
  \ifx\insertframesubtitle\@empty\else
    \vspace*{0.15em}{\usebeamerfont{framesubtitle}\usebeamercolor[fg]{framesubtitle}\insertframesubtitle\par}%
  \fi
  \vspace*{0.25em}{\color{accent}\rule{\textwidth}{0.6pt}}\par}
\makeatother
\setbeamertemplate{footline}{%
  \hspace*{1.4em}{\tiny\color{muted}\eventline}\hfill{\tiny\color{muted}\insertframenumber\,/\,\inserttotalframenumber}\hspace*{1.4em}\vspace*{0.8em}}
\setbeamertemplate{title page}{%
  \begin{tikzpicture}[remember picture, overlay]
    \fill[accent] (current page.north west) rectangle ([xshift=5pt]current page.south west);
  \end{tikzpicture}%
  \begin{columns}[c]
  \begin{column}{0.64\textwidth}
  {\footnotesize\bfseries\color{accent}\MakeUppercase{\insertdate}\par}
  \vspace{1.4em}{\usebeamerfont{title}\usebeamercolor[fg]{title}\inserttitle\par}
  \vspace{0.9em}{\small\insertsubtitle\par}
  \vspace{1.6em}{\footnotesize\insertauthor\par}
  \end{column}
  \begin{column}{0.3\textwidth}
  \centering
  %%TITLEPIC%%

  \vspace{0.3em}{\scriptsize\color{muted}OpenAI's pattern $H$, a $66\times66$-matrix}
  \end{column}
  \end{columns}
}

\usepackage{amsmath, amssymb}
\usepackage{tikz}
\usetikzlibrary{arrows.meta, decorations.markings, decorations.pathreplacing, positioning, calc}
\DeclareMathOperator{\dist}{dist}
\newtheorem{conjecture}[theorem]{Conjecture}
\tikzset{
  tnode/.style={circle, draw, thick, minimum size=8mm, inner sep=0pt, font=\small},
  note/.style={font=\scriptsize, align=center},
  lab/.style={font=\scriptsize},
  forced/.style={line width=2pt, blue!70!black, opacity=0.45, line cap=round,
    postaction={decorate}, decoration={markings,
      mark=at position #1 with {\arrow[opacity=1]{Stealth[length=2.6mm, width=2.4mm]}}}},
  forced/.default=0.75,
  dropped/.style={forced=#1, red!80!black, opacity=0.6},
  dropped/.default=0.75,
}
\newcommand{\treepart}[2]{%
  \foreach \row [count=\i from 0] in {{0,0,0,0,0,0},{0,0,0,0,0,0},{1,1,#1,1,0,0},
                                       {1,1,0,1,0,0},{1,1,0,0,0,0},{1,1,0,0,0,0}} {
    \foreach \v [count=\j from 0] in \row {
      \ifnum\v=1 \fill[gray!25] (\j, -\i) rectangle (\j+1, -\i-1); \fi
    }
  }
  #2
  \foreach \row [count=\i from 0] in {{0,0,0,0,0,0},{0,0,0,0,0,0},{1,1,#1,1,0,0},
                                       {1,1,0,1,0,0},{1,1,0,0,0,0},{1,1,0,0,0,0}} {
    \foreach \v [count=\j from 0] in \row {
      \draw[gray!60] (\j, -\i) rectangle (\j+1, -\i-1);
      \node[font=\small] at (\j+0.5, -\i-0.5) {$\v$};
    }
  }
  \foreach \x in {2,3,4} \draw[thick] (\x, 0) -- (\x, -6);
  \foreach \y in {2,3,4} \draw[thick] (0, -\y) -- (6, -\y);
  \draw[thick] (0, 0) rectangle (6, -6);
  \node[lab, anchor=east] at (-0.1, -1)   {$R^-$};
  \node[lab, anchor=east] at (-0.1, -2.5) {$r_0$};
  \node[lab, anchor=east] at (-0.1, -3.5) {$r_1$};
  \node[lab, anchor=east] at (-0.1, -5)   {$R^+$};
  \node[lab, anchor=south] at (1, 0.05)   {$C^+$};
  \node[lab, anchor=south] at (2.5, 0.05) {$c_0$};
  \node[lab, anchor=south] at (3.5, 0.05) {$c_1$};
  \node[lab, anchor=south] at (5, 0.05)   {$C^-$};
  \draw[ultra thick, blue!70!black] (0, -4) rectangle (2, -6);
  \draw[ultra thick, orange!90!black] (4, 0) rectangle (6, -2);
  \draw[ultra thick, red!80!black] (2.06, -2.06) rectangle (2.94, -2.94);
  \draw[ultra thick, red!80!black] (3.06, -3.06) rectangle (3.94, -3.94);
}
\tikzset{leanbox/.style={draw=ink!60, rounded corners=2pt, fill=white, align=center, inner sep=2.5pt, text width=#1, font=\scriptsize},
  leanbox/.default=2.3cm}
% layout helpers: a callout for the sentence that carries a step, a coloured step heading,
% and a tight list for the text of an overlay step
\newcommand{\keybox}[1]{\par\smallskip\noindent\hbox{{\color{accent}\vrule width 2pt}{\setlength{\fboxsep}{4pt}\colorbox{accent!9}{\parbox{\dimexpr\linewidth-2\fboxsep-2pt\relax}{\raggedright #1}}}}\par}
\newcommand{\stephead}[1]{{\color{accent!85!black}\bfseries #1}\par\vspace{0.15em}}
\newenvironment{steplist}{\setlength{\leftmargini}{1.1em}\begin{itemize}\setlength{\itemsep}{0.18em}\setlength{\parskip}{0pt}\setlength{\topsep}{0pt}}{\end{itemize}}
\newcommand{\smallmat}[1]{\left(\begin{smallmatrix}#1\end{smallmatrix}\right)}

% title, subtitle, authors with affiliations and date as in the team's report
\title{Polynomial removal for binary matrices}
\subtitle{Problem \#96 at the Caltech \texorpdfstring{$\times$}{×} SAIR Mathathon}
\author{\texorpdfstring{%
\textbf{Team costco}\\[0.6em]
\begin{tabular}{@{}l@{\hspace{1.2em}}l@{}}
Papangkorn Apinyanon & \textcolor{muted}{MIT}\\
Saksham Sethi & \textcolor{muted}{Caltech}\\
Henning Arn\'{o}r Skeggi \'{U}lfarsson & \textcolor{muted}{Reykjavik University}
\end{tabular}}{Papangkorn Apinyanon, Saksham Sethi, Henning Arn\'{o}r Skeggi \'{U}lfarsson}}
\date{November 15, 2026}

\begin{document}

\begin{frame}[c]
\titlepage
\end{frame}

% ------------------------------------------------------------------ motivation
\begin{frame}{Why removal?}
\framesubtitle{Far from pattern-free should mean many copies, and then a few random samples find one}
\begin{columns}[T]
\begin{column}{0.58\textwidth}
\small
\begin{itemize}
\item<1-> \textbf{Removal lemmas.} If many entries must change to destroy every copy of a small pattern, there are many copies. A central tool: triangle removal implies Roth's theorem.
\item<2-> \textbf{The engine behind property testing.} Is a huge matrix $A$ free of a pattern $H$, or far from it? Sample rows and columns at random and look: here $%%MOT_N%%$ of the $%%MOT_TOTAL%%$ choices of $2$ rows and $2$ columns read $H$.
\item<3-> \textbf{How many tries?} About $1/\delta$ if $A$ holds a $\delta$ share of all possible copies (here about $%%MOT_TRIES%%$). If $A$ is far from $H$-free, is $\delta$ only \alert{polynomially} small?
\end{itemize}
\end{column}
\begin{column}{0.38\textwidth}
\centering
%%MOTIVPIC%%
\end{column}
\end{columns}
\end{frame}

\begin{frame}{What we set out to do}
\framesubtitle{Make OpenAI's counterexample readable; turn single counterexamples into families}
\small
Polynomial removal for ordered binary matrices was a central open problem for nearly two decades. OpenAI's $66\times66$ counterexample settled it in the negative, and smaller counterexamples followed within weeks, at least one found with the help of AI.

\medskip
\begin{columns}[T]
\begin{column}{0.48\textwidth}
\begin{block}{1. Understand it}
The constructions are a tangle of anchors, modes and recursive parts. Why does OpenAI's counterexample escape the approach of Alon and Ben-Eliezer? We give a human-readable account.
\end{block}
\end{column}
\begin{column}{0.48\textwidth}
\begin{block}{2. Extend it}
Counterexamples were found one pattern at a time. \textbf{Direct sums} and \textbf{duplication} give the first infinite families of counterexamples, and new families that satisfy the conjecture.
\end{block}
\end{column}
\end{columns}

\medskip
\centering
\begin{tikzpicture}[>={Stealth[length=2mm]}, node distance=0.32cm, every node/.append style={minimum height=0.8cm}]
\node[leanbox=2.15cm, fill=accent!8] (a) {copies and\\ the conjecture};
\node[leanbox=2.15cm, fill=accent!8, right=of a] (b) {disjoint copies\\ are not enough};
\node[leanbox=2.15cm, fill=accent!8, right=of b] (c) {OpenAI's\\ construction};
\node[leanbox=2.15cm, fill=teal!10, right=of c] (d) {direct sums and\\ duplication};
\node[leanbox=2.15cm, fill=teal!10, right=of d] (e) {Lean};
\foreach \u/\v in {a/b, b/c, c/d, d/e} \draw[->, muted] (\u) -- (\v);
\end{tikzpicture}
\end{frame}

% ------------------------------------------------------------------ basics
\begin{frame}{Ordered copies}
\framesubtitle{Both $0$s and $1$s must be matched}
\begin{columns}[T]
\begin{column}{0.58\textwidth}
\small
\begin{block}{Definition}
A \textbf{copy} of a $k\times k$ pattern $H$ in an $n\times n$ binary matrix $A$:
rows $r_1<\dots<r_k$ and columns $c_1<\dots<c_k$ with
$A(r_i,c_j)=H(i,j)$ for all $i,j$.
$N_H(A)$ is the number of copies.
\end{block}

{\color{accent!85!black}\bfseries Example.} $H=\smallmat{1&0\\1&1}$ in the matrix on the right:
\begin{steplist}
\item<2-> rows $1,2$, columns $1,3$
\item<3-> rows $1,2$, columns $1,2$
\item<4-> rows $1,3$, columns $1,3$
\end{steplist}
\uncover<5->{\keybox{No others: $N_H(A)=3$, and all three use the entry $A(1,1)$.}}
\end{column}
\begin{column}{0.38\textwidth}
\centering
%%COPIES%%
\end{column}
\end{columns}
\end{frame}

\begin{frame}{Removing copies}
\framesubtitle{Sometimes one entry meets every copy}
\begin{columns}[T]
\begin{column}{0.55\textwidth}
\begin{block}{Definition}
\begin{itemize}
\item $D_H(A)$: the fewest entries of $A$ to change so that no copy of $H$ is left.
\item $\dist_H(A)=D_H(A)/n^2$, and $A$ is \textbf{$\epsilon$-far} from $H$-free if $\dist_H(A)\ge\epsilon$.
\end{itemize}
\end{block}

\medskip
\stephead{Example}
\begin{overprint}
\onslide<1>
\raggedright
All three copies use $A(1,1)$.
\onslide<2->
\raggedright
Change $A(1,1)$ to $0$: no copy is left. One change was needed, so $D_H(A)=1$ and $\dist_H(A)=1/9$.
\end{overprint}
\end{column}
\begin{column}{0.4\textwidth}
\centering
%%DISTANCE%%
\end{column}
\end{columns}
\end{frame}

\begin{frame}{A change can create a new copy}
\framesubtitle{Repairing one copy can make another}
\begin{columns}[T]
\begin{column}{0.55\textwidth}
\begin{itemize}
\item<1-> A copy: rows $1,2$, columns $1,2$.
\item<2-> Repair it by changing $A(1,2)$ from $0$ to $1$.
\item<3-> That copy is gone \dots
\item<4-> \dots\ but rows $1,2$ with columns $2,3$ now read $\smallmat{1&0\\1&1}$: a \alert{new copy}.
\item<5-> Still $3$ copies. A repair has to watch the copies it creates.
\end{itemize}
\end{column}
\begin{column}{0.4\textwidth}
\centering
%%CREATE%%
\end{column}
\end{columns}
\end{frame}

\begin{frame}{The conjecture}
\begin{conjecture}[polynomial removal]
For every $k\times k$ binary matrix $H$ there are $c,C>0$ such that for every $n$, every
$n\times n$ binary matrix $A$ and every $\epsilon\in(0,1]$:
\[ \dist_H(A)\ge\epsilon \quad\Longrightarrow\quad N_H(A)\ge c\,\epsilon^{C}n^{2k}. \]
\end{conjecture}
\medskip
\begin{itemize}
\item If many changes are needed to remove the copies, there should be many copies: at least a
fraction polynomial in $\epsilon$ of the about $n^{2k}$ possible positions.
\item For a $k\times l$ pattern, read the conjecture with $n^{k+l}$ in place of $n^{2k}$.
\end{itemize}
\end{frame}

\begin{frame}{History}
\small
\setbeamercolor{description item}{fg=accent!85!black}
\begin{description}[Nov 2026]\setlength{\itemsep}{0.22em}
\item[2007] Alon, Fischer, Newman: polynomial removal without row and column order; the ordered case is raised as an open direction.
\item[2017] Alon, Ben-Eliezer, Fischer: removal holds for ordered matrices: an $\epsilon$-far matrix has at least $\delta n^{2k}$ copies, for some $\delta>0$ depending on $\epsilon$ and $H$. But the $1/\delta$ their proof gives is of wowzer type in $1/\epsilon$, not polynomial.
\item[2020] Alon, Ben-Eliezer: the question for finite families (Problem~1.4); polynomial removal when there are many \emph{disjoint} copies.
\item[2025] Gishboliner, Shapira, survey: ``one of the key open problems'' (Conjecture~4.5).
\item[Sep 2026] OpenAI's preprint: the conjecture \alert{fails} for a $66\times66$ pattern.
\item[Oct 2026] smaller ones: a $5\times5$ draft (ProveIt, 7 Oct), a $4\times4$ with a Lean proof (Erebos Labs, 9 Oct), the same $4\times4$ on arXiv (Wang, 11 Oct), a claimed $3\times3$ (Shea, 12 Oct, unrefereed).
\item[Nov 2026] Lupini: the $2\times2$ identity \emph{does} have polynomial removal.
\end{description}
\end{frame}

% ------------------------------------------------------------------ Alon and Ben-Eliezer
\begin{frame}{The approach of Alon and Ben-Eliezer}
\framesubtitle{Many disjoint copies give many copies}
\begin{columns}[T]
\begin{column}{0.6\textwidth}
Copies are \textbf{pairwise disjoint} if no two share an entry.

\begin{block}{Theorem (Alon--Ben-Eliezer)}
If $A$ has $\epsilon n^2$ pairwise disjoint copies of $H$,
then $N_H(A)\ge\delta n^{2k}$ with $1/\delta$ polynomial in $1/\epsilon$.
\end{block}

So it would suffice to show: \emph{$\epsilon$-far $\Rightarrow$ $\tau n^2$ pairwise disjoint copies},
$1/\tau$ polynomial in $1/\epsilon$.

\uncover<3->{In the example all copies share one entry: at most one disjoint copy, but then one
change removes them all.
\keybox{\alert{Can few disjoint copies need many changes?}}}
\end{column}
\begin{column}{0.35\textwidth}
\centering
%%DISJOINT%%
\end{column}
\end{columns}
\end{frame}

\begin{frame}{How OpenAI's hosts $A_h$ escape}
\framesubtitle{Few disjoint copies, yet far from $H$-free}
For OpenAI's $66\times66$ pattern $H$ and its hosts $A_h$, with $m=2^h$ and $n=(386h+2)m$:
\begin{enumerate}
\item<2-> $A_h$ is $\epsilon$-far from $H$-free for $\epsilon=(386h+2)^{-2}$: every repair changes at least $m^2=\epsilon n^2$ entries \hfill (OpenAI's Prop.~3.1)
\item<3-> every copy of $H$ passes through one of only $m$ \textbf{leaf cells}, so there are at most $m$ pairwise disjoint copies \hfill (OpenAI's Prop.~4.6)
\item<4-> so $\tau\le m/n^2=\epsilon/m$, and $m=2^{(\epsilon^{-1/2}-2)/386}$ is exponential in $\epsilon^{-1/2}$.
\end{enumerate}
\medskip
\uncover<5->{The conjecture itself fails too: $A_h$ has at most $\epsilon\,2^{-h}n^{132}$ copies, smaller than
any fixed power of $\epsilon$ times $n^{132}$ as $h$ grows.}

\medskip
\uncover<6->{\keybox{\alert{Few copies, all through $m$ entries, yet $m^2$ changes needed.} How?}}
\end{frame}

% ------------------------------------------------------------------ the construction
\begin{frame}{OpenAI's pattern $H$}
\framesubtitle{$64$ of its $66$ rows and columns are an address; the content is a $2\times2$ body}
\begin{columns}[c]
\begin{column}{0.48\textwidth}
\[
H=\begin{pmatrix} S & e_1 & e_2\\ e_1^{\mathsf T} & 1 & 0\\ e_2^{\mathsf T} & 1 & 1\end{pmatrix}
\]
\begin{itemize}
\item \textbf{anchor} $S$, $64\times64$: $1$ off the diagonal, $0$ on it, except that rows $33$--$64$ spell the $32$ five-bit words in columns $60$--$64$
\item \textbf{marks} $e_1,e_2$: point back to the first two rows and columns of $S$
\item \textbf{body} $P=\smallmat{1&0\\1&1}$
\end{itemize}
\end{column}
\begin{column}{0.48\textwidth}
\centering
%%HPICTURE%%
\end{column}
\end{columns}
\end{frame}

\begin{frame}{The host $A_h$}
\framesubtitle{Anchors, a tree part and dummies}
\begin{columns}[c]
\begin{column}{0.5\textwidth}
\begin{itemize}
\item \textbf{anchor part}: $6h$ blown-up copies of $S$ (each row and column repeated $m$ times), one for each \textbf{mode}, as diagonal blocks that meet only in $0$s
\item \textbf{tree part}: built on a complete binary tree of depth $h$ with $m=2^h$ leaves; every node owns rows and columns; a leaf's row and column cross in its \textbf{leaf cell}
\item \textbf{dummy} rows and columns: all $1$s
\item size $n=(386h+2)m$
\end{itemize}
\end{column}
\begin{column}{0.46\textwidth}
\centering
\begin{tikzpicture}[x=0.85cm,y=0.85cm,font=\scriptsize]
\draw[thick] (0,0) rectangle (5,-5);
\foreach \k in {0,...,5} { \fill[gray!35] (\k*0.5,-\k*0.5) rectangle (\k*0.5+0.5,-\k*0.5-0.5); }
\draw (3,0) -- (3,-5); \draw (0,-3) -- (5,-3);
\fill[teal!15] (3,0) rectangle (5,-3); \fill[teal!15] (0,-3) rectangle (3,-5);
\fill[orange!18] (3,-3) rectangle (5,-5);
\fill[gray!25] (3,-3) rectangle (3.35,-5); \fill[gray!25] (3,-4.65) rectangle (5,-5);
\foreach \k in {0,...,5} { \fill[gray!45] (\k*0.5,-\k*0.5) rectangle (\k*0.5+0.5,-\k*0.5-0.5); }
\draw[thick] (0,0) rectangle (5,-5); \draw (3,0) -- (3,-5); \draw (0,-3) -- (5,-3);
\node[align=left, anchor=west] at (1.25,-0.6) {anchors: one\\block per mode};
\node[align=center] at (4,-1.5) {marks};
\node[align=center] at (1.5,-4) {marks};
\node[align=center] at (4.15,-3.9) {tree\\part};
\node[anchor=west] at (5.1,-4.82) {dummy row};
\node[anchor=north] at (3.17,-5.05) {dummy column};
\end{tikzpicture}
\end{column}
\end{columns}
\end{frame}

\begin{frame}{Where a copy of $H$ can sit}
\framesubtitle{In $A_1$: one mode's anchor, its marks, and a $2\times2$ body}
\begin{columns}[T]
\begin{column}{0.52\textwidth}
\centering
%%HOSTVIEW%%
\end{column}
\begin{column}{0.45\textwidth}
\small
\begin{overprint}
\onslide<1>
\raggedright
\stephead{All of $A_1$}
\begin{steplist}
\item Each anchor shrunk: groups $1$ and $2$ get one row and column each, groups $3$--$64$ one small block (drawn as $S$).
\item Six modes, six anchors on the diagonal, $0$s between them.
\item Bottom right, drawn larger: the tree part and the dummies, one cell per row and column.
\end{steplist}
\onslide<2>
\raggedright
\stephead{The anchor}
\begin{steplist}
\item A copy of $S$ can only use one mode's block, with one row and one column from each of its $64$ groups ($2^{128}$ ways).
\item Here the mode is $W_{1,0}^-$.
\end{steplist}
\onslide<3>
\raggedright
\stephead{The marks}
\begin{steplist}
\item In this mode's anchor, row group $1$ has its $1$s in the mode's \textcolor{teal!80!black}{first} columns, row group $2$ in its \textcolor{orange!70!yellow!90!black}{second} columns.
\item The column groups likewise point to the first and second rows (ticks on the right and below).
\end{steplist}
\onslide<4>
\raggedright
\stephead{The body}
\begin{steplist}
\item So the last two rows of $H$ must be a first and a second row ($r_0$ and a dummy row),
\item its last two columns a first and a second column ($c_0$ and a column of $C_0^-$).
\end{steplist}
\keybox{They cross in a $2\times2$ block, which must read $P=\smallmat{1&0\\1&1}$: here it does.}
\onslide<5>
\raggedright
\stephead{A copy of $H$}
These $66$ rows and $66$ columns (shaded) read exactly $H$.
\onslide<6>
\raggedright
\stephead{All copies}
\begin{steplist}
\item $A_1$ has %%A1_BODIES%% bodies, all at the two leaf cells (circled), in the modes $W_{1,0}^-$ and $W_{1,1}^-$.
\end{steplist}
\keybox{With the $2^{128}$ anchor choices each, they give every copy of $H$ (OpenAI's Prop.~4.6).}
\end{overprint}
\end{column}
\end{columns}
\end{frame}

\begin{frame}{Why $S$ cannot hide in the tree part}
\framesubtitle{The tree's rows are too regular to show all $32$ five-bit words}
\begin{columns}[T]
\begin{column}{0.52\textwidth}
\centering
%%WORDSVIEW%%
\end{column}
\begin{column}{0.45\textwidth}
\small
\begin{overprint}
\onslide<1>
\raggedright
Rows $33$--$64$ of $S$, read on its columns $60$--$64$, are all $32$ words of five bits.
\onslide<2>
\raggedright
The tree and dummy part of $A_2$, read on five of its columns: the five that show the most words.
\onslide<3->
\raggedright
\keybox{Only \alert{%%WORDS_FOUND%% of the $32$} words appear; for every $h$ at most $27$ do (OpenAI's Lemma~4.3).}
\begin{steplist}
\item So a copy of $S$ cannot have its rows $33$--$64$ and its columns $60$--$64$ in this part: more than $32$ of its rows or columns are anchors.
\item The dense half of $S$ then keeps the copy inside one mode's block (OpenAI's Lemma~4.4).
\end{steplist}
\end{overprint}
\end{column}
\end{columns}
\end{frame}

\begin{frame}{Each body is a forced step, and the steps form chains}
\framesubtitle{A chain runs from $1$s to $0$s, so it must drop somewhere, and the drop is a copy}
\begin{columns}[T]
\begin{column}{0.49\textwidth}
\small
\begin{minipage}[t]{0.28\linewidth}
\vspace{0pt}
\centering
\begin{tikzpicture}[x=0.62cm,y=0.62cm]
\fill[gray!22] (0,0) rectangle (1,-2);
\draw[gray!60] (1,0) -- (1,-2); \draw[gray!60] (0,-1) -- (2,-1); \draw (0,0) rectangle (2,-2);
\node at (0.5,-0.5) {$1$}; \node at (0.5,-1.5) {$1$};
\node at (1.5,-0.5) {$x$}; \node at (1.5,-1.5) {$y$};
\node[font=\scriptsize, anchor=south] at (0.5,0.02) {dummy};
\draw[-{Stealth[length=2.2mm]}, accent!90!black, very thick] (1.5,-1.25) -- (1.5,-0.75);
\end{tikzpicture}
\end{minipage}\hfill
\begin{minipage}[t]{0.68\linewidth}
\vspace{0pt}
\raggedright
One row or column of every body is a \textbf{dummy}, all $1$s. The body reads $P=\smallmat{1&0\\1&1}$ exactly when $y=1$ and $x=0$.
\end{minipage}

\keybox{\alert{To avoid $H$, a $1$ at $y$ forces a $1$ at $x$}: a \textbf{forced step}.}
\vspace{-0.3em}
\begin{overprint}
\onslide<1>
\raggedright
\onslide<2->
\raggedright
\begin{steplist}
\item The modes link forced steps into \textbf{chains}, from the root's block of $1$s through a leaf cell to its block of $0$s.
\item Every chain must \alert{drop} from $1$ to $0$ somewhere, and each drop gives copies of $H$.
\item In $A_1$ each chain drops right after its leaf cell.
\end{steplist}
\end{overprint}
\end{column}
\begin{column}{0.49\textwidth}
\centering
\begin{tikzpicture}[baseline=(root.north), tnode/.append style={minimum size=6.5mm, font=\scriptsize}]
\node[tnode] (root) at (0, 0) {root};
\node[tnode] (l0) at (-0.65, -1.3) {$0$};
\node[tnode] (l1) at ( 0.65, -1.3) {$1$};
\draw[thick] (root) -- (l0);
\draw[thick] (root) -- (l1);
\node[note, below=1mm of l0] {$r_0, c_0$};
\node[note, below=1mm of l1] {$r_1, c_1$};
\node[note, above=1mm of root] {$R^\pm, C^\pm$};
\end{tikzpicture}\hspace{0.2cm}%
\begin{tikzpicture}[x=5.4mm, y=5.4mm, baseline=(top)]
\coordinate (top) at (0, 0.6);
\only<1>{\treepart{1}{}}
\only<2->{\treepart{1}{%
  \draw[forced]  (0.5, -4.5) -- (0.5, -2.5);
  \draw[forced]  (0.5, -2.5) -- (2.5, -2.5);
  \draw[dropped] (2.5, -2.5) -- (4.5, -2.5);
  \draw[forced=0.5] (4.5, -2.5) -- (4.5, -1.5);
  \draw[forced]  (1.5, -5.5) -- (1.5, -3.5);
  \draw[forced]  (1.5, -3.5) -- (3.5, -3.5);
  \draw[dropped] (3.5, -3.5) -- (5.5, -3.5);
  \draw[forced]  (5.5, -3.5) -- (5.5, -1.5);
}}
\end{tikzpicture}

\vspace{0.3em}
{\scriptsize\color{muted} The rows and columns of $A_1$ that the tree's nodes own. Blue box: the root's block of $1$s; orange box: its block of $0$s; red: the leaf cells.}
\end{column}
\end{columns}
\end{frame}

\begin{frame}{Hitting every copy is cheap, removing them is not}
\framesubtitle{In $A_1$: repairing one entry makes another copy appear}
\begin{columns}[T]
\begin{column}{0.42\textwidth}
\centering
%%CORE%%
\end{column}
\begin{column}{0.55\textwidth}
\small
{\scriptsize\color{muted} The tree and dummy part of $A_1$ ($m=2$). Orange arrows: drops, with their number of bodies; green: changed entries; circles: leaf cells. Each body gives $2^{128}$ copies of $H$.\par}

\begin{overprint}
\onslide<1>
\raggedright
\stephead{$A_1$}
\begin{steplist}
\item Every chain drops right after its leaf cell: \alert{%%CORE_BODIES_1%% bodies}, all through the $2$ leaf cells.
\end{steplist}
\onslide<2>
\raggedright
\stephead{Change the leaf cell $(r_0,c_0)$ (%%CORE_CHANGED_2%% change)}
\begin{steplist}
\item Its copies are gone, but the step into it now drops:
\item \alert{%%CORE_BODIES_2%% bodies}, as many as before.
\end{steplist}
\onslide<3>
\raggedright
\stephead{Change both leaf cells (%%CORE_CHANGED_3%% changes)}
\begin{steplist}
\item Every original copy is hit,
\item yet every chain just drops one step earlier: \alert{%%CORE_BODIES_3%% bodies} (OpenAI's Remark~4.7).
\end{steplist}
\onslide<4>
\raggedright
\stephead{Cut the chains instead}
\begin{steplist}
\item Set the dummy rows to $0$ in both leaf columns, $m^2=%%CORE_CHANGED_4%%$ changes.
\item \alert{%%CORE_BODIES_4%% bodies}: no copy of $H$ is left (checked by computer).
\end{steplist}
\onslide<5>
\raggedright
\stephead{No repair is cheaper (OpenAI's Prop.~3.1)}
\begin{steplist}
\item A repair must break every chain, at one of its ends, its dummy entries or its anchor entries,
\item and each such entry lies on at most a $1/m^2$ share of the chains, while a leaf cell lies on a $1/m$ share.
\end{steplist}
\keybox{So hitting every copy costs $m$ changes; removing them costs $m^2=\epsilon n^2$.}
\end{overprint}
\end{column}
\end{columns}
\end{frame}

% ------------------------------------------------------------------ our results
\begin{frame}{Our results: two reductions}
\begin{block}{Direct sums}
If $H$ is a counterexample with no all-zero row or column, then so is
$\smallmat{H&0\\0&Z}$ for every $Z$ with pairwise distinct rows or pairwise distinct columns.
\end{block}
\begin{block}{Duplication}
Repeating a row or a column next to itself does not change whether a pattern satisfies the conjecture.
\end{block}
\medskip
\keybox{Consequences: the first \textbf{infinite families of counterexamples}, and, with Lupini's theorem
for the $2\times2$ identity, an infinite family of patterns that \textbf{satisfy} the conjecture.}
\end{frame}

\begin{frame}{Direct sums: the construction}
\framesubtitle{Attach $Z$ diagonally, and blow it up in the host}
\begin{columns}[c]
\begin{column}{0.56\textwidth}
\small
$H$ is $k\times l$, $Z$ is $p\times q$, $A$ is $n\times n$:
\[
H'=\begin{pmatrix} H & 0\\ 0 & Z\end{pmatrix},
\qquad
A^+=\begin{pmatrix} A & 0\\ 0 & \widehat Z\end{pmatrix},
\]
where $\widehat Z$ replaces each entry of $Z$ by a constant $qn\times pn$ block; $A^+$ has size $m=(1+pq)n$.

\begin{block}{Theorem}
If $H$ has no all-zero row or column, and $Z$ has distinct rows or distinct columns, then
\[ D_{H'}(A^+)=D_H(A),\qquad N_{H'}(A^+)\le m^{p+q}N_H(A). \]
\end{block}
\end{column}
\begin{column}{0.4\textwidth}
\centering
\begin{tikzpicture}[x=1cm,y=1cm]
\fill[blue!10] (0,0) rectangle (2.2,-2.2);
\fill[gray!20] (2.2,-2.2) rectangle (4,-4);
\draw[thick] (0,0) rectangle (4,-4);
\draw (2.2,0) -- (2.2,-4); \draw (0,-2.2) -- (4,-2.2);
\node at (1.1,-1.1) {$A$};
\node at (3.1,-1.1) {$0$};
\node at (1.1,-3.1) {$0$};
\node at (3.1,-3.1) {$\widehat Z$};
\node[font=\scriptsize, anchor=south] at (1.1,0.05) {old};
\node[font=\scriptsize, anchor=south] at (3.1,0.05) {new};
\end{tikzpicture}
\end{column}
\end{columns}
\end{frame}

\begin{frame}{Direct sums: an example}
\framesubtitle{The smallest case, $Z=(1)$: one new row and one new column}
\begin{columns}[T]
\begin{column}{0.55\textwidth}
\small
Take $H=\smallmat{1&0\\1&1}$ ($k=l=2$), $Z=(1)$ ($p=q=1$), and as host $A$ the $3\times3$ example ($n=3$). Then $H'=\smallmat{H&0\\0&Z}=%%DS_HPRIME%%$.
\begin{overprint}
\onslide<1>
\raggedright
\stephead{The host}
\begin{steplist}
\item $\widehat Z$ replaces the entry $1$ of $Z$ by a constant block with $qn=%%DS_NEW%%$ rows and $pn=%%DS_NEW%%$ columns.
\item So $A^+$ has size $m=(1+pq)n=%%DS_M%%$: rows and columns %%DS_OLD%% are old, %%DS_NEWR%% new.
\item The new rows form one group of identical rows, and the new columns one group of identical columns.
\end{steplist}
\onslide<2>
\raggedright
\stephead{A copy of $H'$}
\begin{steplist}
\item Its two H-rows and two H-columns are \textcolor{accent!90!black}{rows $%%DS_HROWS%%$ and columns $%%DS_HCOLS%%$}, a copy of $H$ in $A$.
\item Its Z-row and Z-column are the new \textcolor{teal!80!black}{row $%%DS_ZROW%%$ and column $%%DS_ZCOL%%$}.
\item Between the two parts $A^+$ has $0$s, as $H'$ needs.
\end{steplist}
\onslide<3>
\raggedright
\stephead{All copies}
\begin{steplist}
\item The H-part of every copy lies in the old part (next slide).
\item So a copy of $H'$ is one of the $N_H(A)=%%DS_NH%%$ copies of $H$ in $A$, plus any of the $%%DS_NEW%%$ new rows and any of the $%%DS_NEW%%$ new columns:
\end{steplist}
\keybox{$N_{H'}(A^+)=%%DS_NH%%\cdot%%DS_NEW%%\cdot%%DS_NEW%%=%%DS_NH2%%$, within the theorem's bound $m^{p+q}N_H(A)=%%DS_BOUND%%$.}
\onslide<4>
\raggedright
\stephead{Distance}
\begin{steplist}
\item Changing $A(1,1)$ removes every copy of $H$ from $A$, so $D_H(A)=%%DS_DH%%$.
\item The same change in $A^+$ removes every copy of $H'$, so
\end{steplist}
\keybox{\centering $D_{H'}(A^+)=%%DS_DH%%=D_H(A)$.}
{\footnotesize (This $H$ only illustrates the theorem; whether it has polynomial removal is open.)}
\end{overprint}
\end{column}
\begin{column}{0.42\textwidth}
\centering
%%DIRECTSUM%%
\end{column}
\end{columns}
\end{frame}

\begin{frame}{Direct sums, proof (1): the $H$-part of every copy is old}
\framesubtitle{New rows have $1$s only in new columns, and they come in too few groups}
\begin{columns}[T]
\begin{column}{0.55\textwidth}
\footnotesize
Let $M=\smallmat{X&0\\0&\widehat Z}$ with $X$ any $n\times n$ matrix. Take a copy of $H'$ in $M$ on rows $r_1<\dots<r_{k+p}$ and columns $c_1<\dots<c_{l+q}$: first its \textbf{H-rows} $r_1,\dots,r_k$ and \textbf{H-columns} $c_1,\dots,c_l$, then its \textbf{Z-rows} and \textbf{Z-columns}.

\keybox{\textbf{Claim.} The H-rows and H-columns are old, so they form a copy of $H$ in $X$.}

\begin{overprint}
\onslide<1>
\raggedright
\onslide<2>
\raggedright
\begin{steplist}
\item Suppose the last H-row $r_k$ is \alert{new}.
\item The Z-rows $r_{k+1},\dots,r_{k+p}$ come after it, so they are new too.
\end{steplist}
\onslide<3>
\raggedright
\begin{steplist}
\item Row $k$ of $H$ contains a $1$, say in column $j\le l$, so $M(r_k,c_j)=1$.
\item A new row has $1$s only in new columns, so $c_j$ is new, and so are the Z-columns $c_{l+1},\dots,c_{l+q}$ after it.
\end{steplist}
\onslide<4>
\raggedright
\begin{steplist}
\item These $p+1$ rows and $q+1$ columns lie in $\widehat Z$ and must read $\smallmat{1&0\\0&Z}$, whose rows are pairwise distinct when those of $Z$ are.
\item But the new rows form only $p$ groups of identical rows, so two of the $p+1$ rows are equal: a \alert{contradiction}.
\item Here $r_k$ and $r_{k+1}$ both lie in group~1 and read $%%DSP_READ_K%%$, where $r_k$ needs $%%DSP_NEED_K%%$ and $r_{k+1}$ needs $%%DSP_NEED_K1%%$.
\end{steplist}
\onslide<5>
\raggedright
\begin{steplist}
\item So $r_k$ is old, and with it $r_1,\dots,r_k$.
\item Column $l$ of $H$ contains a $1$, say in row $i\le k$, so $M(r_i,c_l)=1$; an old row has $1$s only in old columns, so $c_l$ is old, and with it $c_1,\dots,c_l$.
\end{steplist}

{\footnotesize If instead the columns of $Z$ are pairwise distinct, run the same argument on $q+1$ columns and $q$ groups.}
\end{overprint}
\end{column}
\begin{column}{0.43\textwidth}
\centering
%%DSPROOF1%%

{\scriptsize\color{muted} Drawn with $p=q=2$, $Z=%%DSP_Z%%$; shaded: $1$s.}
\end{column}
\end{columns}
\end{frame}

\begin{frame}{Direct sums, proof (2): the count, and $D_{H'}(A^+)\le D_H(A)$}
\framesubtitle{Both follow from (1) with $X=A$ and with $X=A_0$}
\begin{columns}[T]
\begin{column}{0.55\textwidth}
\small
\begin{itemize}
\item<1-> \textbf{The count.} By (1) with $X=A$, every copy of $H'$ in $A^+$ is a copy of $H$ in $A$ together with $p$ more rows and $q$ more columns. There are at most $m^{p+q}$ ways to choose them, so $N_{H'}(A^+)\le m^{p+q}N_H(A)$.
\item<2-> \textbf{A repair.} Change $D_H(A)$ entries of $A$ to get an $H$-free matrix $A_0$, and leave the rest of $A^+$ as it is. By (1) with $X=A_0$, a copy of $H'$ would contain a copy of $H$ in $A_0$, and there is none. So $D_{H'}(A^+)\le D_H(A)$.
\item<3-> \textbf{The other direction is the real work:} a repair of $A^+$ may also change the $0$ blocks and $\widehat Z$, and could be cheaper there.
\end{itemize}
\end{column}
\begin{column}{0.43\textwidth}
\centering
%%DSPROOF2%%

{\scriptsize\color{muted} Orange: a copy of $H$ in $A$. Teal: Z-rows and Z-columns, any new lines.}
\end{column}
\end{columns}
\end{frame}

\begin{frame}{Direct sums, proof (3): no repair of $A^+$ is cheaper}
\framesubtitle{Pick one new line per group; a good pick turns any repair of $A^+$ into a repair of $A$}
\begin{columns}[T]
\begin{column}{0.55\textwidth}
\footnotesize
Let $B$ be any $m\times m$ matrix with no copy of $H'$. $E_0,E_1,E_2$: the numbers of entries where $B$ differs from $A^+$ in the old--old block, the two mixed blocks, and the new--new block (\textcolor{red!75!black}{crosses}). \keybox{\textbf{Goal:} $E_0+E_1+E_2\ge D$, where $D=D_H(A)\le n^2$.}

\begin{overprint}
\onslide<1>
\raggedright
\onslide<2>
\raggedright
\begin{steplist}
\item \textbf{Pick} at random one row $\rho_i$ from each of the $p$ groups of new rows and one column $\gamma_j$ from each of the $q$ groups of new columns.
\item The pick is \textbf{good} if $B(\rho_i,\gamma_j)=Z(i,j)$ for all $i,j$.
\item A fixed new--new entry is picked with probability $\frac1{qn\cdot pn}$, so $\Pr[\text{not good}]\le \frac{E_2}{pq\,n^2}$.
\end{steplist}
\onslide<3>
\raggedright
\begin{steplist}
\item An old row is \textbf{bad} if $B$ has a $1$ where it meets a picked column; an old column, where it meets a picked row.
\item Each such $1$ is one of the $E_1$ changes, and lies on a picked line with probability at most $\frac1n$: on average at most $\frac{E_1}n$ bad lines.
\item $B^\circ$: the old--old block of $B$ with the bad lines set to $0$. It differs from $A$ in at most $E_0+W$ entries, $W=n\cdot\#\text{bad lines}$, $\mathbb E[W]\le E_1$.
\end{steplist}
\onslide<4>
\raggedright
\stephead{Claim: if the pick is good, $B^\circ$ is $H$-free.}
\begin{steplist}
\item Every row and column of $H$ has a $1$, so a copy of $H$ in $B^\circ$ avoids the zeroed lines;
\item its lines are not bad, so they meet the picked lines in $0$s. The picked lines read $Z$.
\end{steplist}
\keybox{Together they form a \alert{copy of $H'$ in $B$}, which is impossible.}
\onslide<5>
\raggedright
\begin{steplist}
\item So a good pick gives $D\le E_0+W$. Taking expectations,
\[ D\Bigl(1-\frac{E_2}{pq\,n^2}\Bigr)\le D\cdot\Pr[\text{good}]\le E_0+E_1, \]
\item and since $D\le n^2\le pq\,n^2$, this gives $E_0+E_1+E_2\ge D$.
\end{steplist}
\keybox{So $D_{H'}(A^+)\ge D_H(A)$.}
\end{overprint}
\end{column}
\begin{column}{0.43\textwidth}
\centering
%%DSPROOF3%%

{\scriptsize\color{muted} $p=q=2$. Teal: picked lines; orange: a copy of $H$ in $B^\circ$.}
\end{column}
\end{columns}
\end{frame}

\begin{frame}{Direct sums: new counterexamples}
\framesubtitle{A host for $H$ with few copies gives one for $H'$}
\small
\keybox{\textbf{Corollary.} If $H$ (no all-zero row or column) is a counterexample, so is $H'=\smallmat{H&0\\0&Z}$ for every $Z$ with pairwise distinct rows or pairwise distinct columns.}
\begin{itemize}
\item<1-> Given $c',C'>0$, put $s=1+pq$, $C=C'$ and $c=c'\,s^{k+l-2C}$. As $H$ is a counterexample, some $n\times n$ matrix $A$ and $\epsilon\in(0,1]$ have $\dist_H(A)\ge\epsilon$ and $N_H(A)<c\,\epsilon^{C}n^{k+l}$.
\item<2-> Take the host $A^+$, of size $m=sn$, and $\epsilon'=\epsilon/s^2$. By the theorem, $\dist_{H'}(A^+)=D_H(A)/(sn)^2\ge\epsilon/s^2=\epsilon'$, and, with $n=m/s$ and $\epsilon=s^2\epsilon'$,
\[ N_{H'}(A^+)\le m^{p+q}N_H(A)<c\,\epsilon^{C}n^{k+l}m^{p+q}=c'\,(\epsilon')^{C'}m^{(k+p)+(l+q)}. \]
\item<3-> So $c'$ and $C'$ fail for $H'$ on $A^+$; as they were arbitrary, $H'$ is a counterexample.
\end{itemize}
\uncover<4->{\textbf{Variant.} $H$ with $q$ all-zero columns appended, with host $\smallmat{A&0\\0&0}$ of size $(1+q)n$: the same holds, by a simpler proof.}
\end{frame}

\begin{frame}{Duplication: repeated rows and columns do not matter}
\framesubtitle{Delete repeats until none are left}
\begin{columns}[T]
\begin{column}{0.6\textwidth}
\begin{block}{Definition}
$\operatorname{red}(H)$: delete every row equal to the row above it and every column equal to the column to its left, until none is left. If $H$ is $k\times l$, $\operatorname{red}(H)$ is $k_0\times l_0$.
\end{block}

\begin{overprint}
\onslide<1>
\raggedright
Row $2$ equals row $1$ \dots
\onslide<2>
\raggedright
\dots\ delete it.
\onslide<3>
\raggedright
Column $2$ equals column $1$ \dots
\onslide<4>
\raggedright
\dots\ delete it: $\operatorname{red}(H)=I_2$.
\end{overprint}

\uncover<4->{\keybox{\textbf{Theorem.} $H$ satisfies the conjecture if and only if $\operatorname{red}(H)$ does.}}

\medskip
\uncover<4->{\footnotesize The proof rests on two facts: Fact 1 counts copies with one more repeat, Fact 2 blows up a host.}
\end{column}
\begin{column}{0.35\textwidth}
\centering
%%RED%%
\end{column}
\end{columns}
\end{frame}

\begin{frame}{Duplication, Fact 1: frames and completions}
\framesubtitle{Copies of $H$ are frames with one more column; two completions make a copy of $H'$}
\begin{columns}[T]
\begin{column}{0.62\textwidth}
\footnotesize
\keybox{\textbf{Fact 1.} Let $H$ be $k\times l$ and $H'$ be $H$ with column $j$ written twice in a row. For every $n\times n$ matrix $B$, with $N=N_H(B)$ and $N'=N_{H'}(B)$: $\;N'\ge \frac{N^2}{2n^{k+l-1}}-\frac N2$.
(For a repeated row, transpose $H$ and $B$.)}

{\color{accent!85!black}\bfseries Example} (right): $H=\smallmat{1&0\\1&1}$, $j=2$, $H'=%%F1_HD%%$.

\begin{overprint}
\onslide<1>
\raggedright
\begin{steplist}
\item A \textbf{frame} $F$: rows $r_1<\dots<r_k$ and columns $c_1<\dots<c_{j-1}<c_{j+1}<\dots<c_l$ of $B$, a would-be copy of $H$ missing its $j$th column.
\item There are $T\le n^{k+l-1}$ frames (here $T=%%F1_T%%$).
\item A column $c$ with $c_{j-1}<c<c_{j+1}$ is a \textbf{completion} of $F$ if adding it gives a copy of $H$; $x_F$ is the number of completions.
\end{steplist}
\onslide<2>
\raggedright
The frame \textcolor{accent!90!black}{rows $%%F1_A_ROWS%%$ with column $%%F1_A_COL%%$} has the completions \textcolor{teal!80!black}{columns %%F1_A_DONE%%}: $x_F=%%F1_A_X%%$.

\stephead{Step 2}
\begin{steplist}
\item Deleting column $j$ of a copy of $H$ leaves a frame, and the deleted column completes it.
\item So the copies of $H$ match the pairs ($F$, a completion of $F$) one to one, and $\sum_F x_F=N$.
\end{steplist}
\onslide<3>
\raggedright
The frame \textcolor{accent!90!black}{rows $%%F1_B_ROWS%%$ with column $%%F1_B_COL%%$} has one completion, \textcolor{teal!80!black}{column $%%F1_B_DONE%%$}. The other $%%F1_REST%%$ frames have none: $\sum_F x_F=%%F1_XSUM%%=%%F1_N%%=N$.
\onslide<4>
\raggedright
\stephead{Step 3}
\begin{steplist}
\item Two completions $c<c'$ of one frame give a \alert{copy of $H'$}: the columns $c_1<\dots<c_{j-1}<c<c'<c_{j+1}<\dots<c_l$ are in order, and $c$, $c'$ both match column $j$.
\item Different ($F$, $\{c,c'\}$) give different copies, so $N'\ge\sum_F\binom{x_F}{2}$.
\item Here: rows $%%F1_COPY_ROWS%%$ with columns $%%F1_COPY_COLS%%$, and $%%F1_BINOMS%%=%%F1_ND%%=N'$.
\end{steplist}
\end{overprint}
\end{column}
\begin{column}{0.34\textwidth}
\centering
%%F1PIC%%
\end{column}
\end{columns}
\end{frame}

\begin{frame}{Duplication, Fact 1: Cauchy--Schwarz}
\framesubtitle{The pairs are fewest when the completions are spread evenly over the frames}
\begin{columns}[T]
\begin{column}{0.56\textwidth}
\small
So far: $\sum_F x_F=N$ over $T\le n^{k+l-1}$ frames, and
\[ N'\ \ge\ \sum_F\binom{x_F}{2}=\frac12\sum_F x_F^2-\frac N2 . \]
\begin{overprint}
\onslide<1>
\raggedright
\stephead{Step 4}
If $N=0$ there is nothing to show. Otherwise, by Cauchy--Schwarz,
\[ N^2=\Bigl(\sum_F 1\cdot x_F\Bigr)^2\le T\sum_F x_F^2 . \]
This is Jensen's inequality for $x\mapsto x^2$: the sum of squares is smallest when every $x_F$ equals the average $N/T$ (right).
\onslide<2>
\raggedright
\stephead{Combining}
\keybox{\centering $N'\ \ge\ \dfrac{N^2}{2T}-\dfrac N2\ \ge\ \dfrac{N^2}{2n^{k+l-1}}-\dfrac N2 . \qquad\Box$}
On the example: $N=%%F1_N%%$, $T=%%F1_T%%$, $\sum_F x_F^2=%%F1_SQ%%$, and indeed $N^2\le T\sum_F x_F^2$.
\onslide<3>
\raggedright
\stephead{How it is used}
If $N_H(B)\ge\delta n^{k+l}$ and $n\ge2/\delta$, then
\[ N_{H'}(B)\ \ge\ \frac{\delta^2}2n^{k+l+1}-\frac\delta2n^{k+l}\ \ge\ \frac{\delta^2}4n^{k+l+1}. \]
\keybox{One repeat turns a copy density $\delta$ into at least $\delta^2/4$, for $H'$ with its $k+l+1$ lines, once $n$ is not too small.}
\end{overprint}
\end{column}
\begin{column}{0.4\textwidth}
\centering
%%F1BARS%%

{\scriptsize\color{muted} Each bar is a frame, its height the number $x_F$ of its completions; a frame gives $\binom{x_F}2$ pairs.}
\end{column}
\end{columns}
\end{frame}

\begin{frame}{Duplication, Fact 2(a): blowing up a host keeps few copies few}
\framesubtitle{A copy of $H$ in the blow-up comes from a copy of $\operatorname{red}(H)$ in the host}
\begin{columns}[T]
\begin{column}{0.57\textwidth}
\footnotesize
\textbf{Blow-up.} $A^{(q)}$: every entry of the $n\times n$ matrix $A$ becomes a constant $q\times q$ block; $A^{(q)}$ is $qn\times qn$.

\keybox{\textbf{Fact 2.} For $H$ ($k\times l$) and $R=\operatorname{red}(H)$ ($k_0\times l_0$):\\
(a) $N_H(A^{(q)})/(qn)^{k+l}\le N_R(A)/n^{k_0+l_0}$;\\
(b) $D_H(A^{(q)})\ge\tfrac12q^2D_R(A)$ for all large enough $q$.}

{\color{accent!85!black}\bfseries Example} (right): $q=2$, $H=\smallmat{1&0\\1&0\\1&1}$, so $R=\smallmat{1&0\\1&1}$.
\begin{overprint}
\onslide<1>
\raggedright
\onslide<2>
\raggedright
\stephead{Proof of (a)}
Take a copy of $H$ in $A^{(q)}$: here \textcolor{accent!90!black}{rows $%%F2A_ROWS%%$ with columns $%%F2A_COLS%%$}.
\onslide<3>
\raggedright
\begin{steplist}
\item Keep one row and one column for each row and column of $R$: \textcolor{teal!80!black}{rows $%%F2A_KROWS%%$, columns $%%F2A_KCOLS%%$}. They read $R$.
\item No two kept rows lie in one block: rows in a block are equal, and $R$ has no two equal neighboring rows (same for columns).
\item So they come from a \textcolor{teal!80!black}{copy of $R$ in $A$}: rows $%%F2A_AROWS%%$, columns $%%F2A_ACOLS%%$.
\end{steplist}
\onslide<4>
\raggedright
\begin{steplist}
\item Each copy of $R$ in $A$ is counted at most $q^{k_0+l_0}(qn)^{k+l-k_0-l_0}$ times: a line in each of its blocks, and the other lines anywhere.
\item Dividing by $(qn)^{k+l}$ gives (a).
\item Here $N_H(A^{(2)})=%%F2A_NH%%\le%%F2A_NR%%\cdot%%F2A_PER%%=%%F2A_BOUND%%$, as $N_R(A)=%%F2A_NR%%$.
\end{steplist}
\end{overprint}
\end{column}
\begin{column}{0.40\textwidth}
\centering
%%F2APIC%%
\end{column}
\end{columns}
\end{frame}

\begin{frame}{Duplication, Fact 2(b): far stays far}
\framesubtitle{Pick one line per block: a cheap repair of $A^{(q)}$ leaves $A$ almost intact}
\begin{columns}[T]
\begin{column}{0.54\textwidth}
\footnotesize
\keybox{\textbf{Claim.} If $q$ is large, changing at most $\tfrac12q^2D_R(A)$ entries of $A^{(q)}$ leaves a copy of $H$. Call the changed matrix $B'$.}
{\color{accent!85!black}\bfseries Example} (right): $q=3$, $D_R(A)=%%F2B_DR%%$, so $%%F2B_BUDGET%%$ changes (green), all in the block of $A(1,1)$, the entry every copy of $R$ in $A$ uses.
\begin{overprint}
\onslide<1>
\raggedright
\onslide<2>
\raggedright
\begin{steplist}
\item \textbf{Pick} one row from each block of rows and one column from each block of columns, at random.
\item They form an $n\times n$ matrix $B$, which is $A$ unless a changed entry is picked.
\item Here \textcolor{teal!80!black}{rows $%%F2B_PR%%$ and columns $%%F2B_PC%%$}: $B=A$.
\end{steplist}
\onslide<3>
\raggedright
\begin{steplist}
\item Each changed entry is picked with probability $1/q^2$, so on average $B$ differs from $A$ in at most $\tfrac12D_R(A)$ entries (here $%%F2B_EXP%%$).
\item By Markov's inequality, with probability at least $\tfrac12$ it differs in fewer than $D_R(A)$, and then $B$ contains $R$.
\item Here $%%F2B_GOOD%%$ of the $%%F2B_TOTAL%%$ picks contain $R$.
\end{steplist}
\onslide<4>
\raggedright
\begin{steplist}
\item A given copy of $R$ in $B'$ is picked with probability at most $q^{-(k_0+l_0)}$, one line from each of its blocks. So
\end{steplist}
\keybox{\centering $N_R(B')\ \ge\ \tfrac12q^{k_0+l_0}$.}
Here $N_R(B')=%%F2B_NRP%%\ge %%F2B_HALF%%$.
\onslide<5>
\raggedright
\begin{steplist}
\item So at least a fraction $\frac{1}{2n^{k_0+l_0}}$ of all positions hold a copy of $R$ in $B'$, and it does not shrink as $q$ grows.
\item Fact 1, once per repeated row or column, gives the same for $H$ once $q$ is large:
\end{steplist}
\keybox{\alert{$B'$ contains $H$}. Here $B'$ has $%%F2B_NHP%%$ copies of $H$.}
\end{overprint}
\end{column}
\begin{column}{0.42\textwidth}
\centering
%%F2BPIC%%
\end{column}
\end{columns}
\end{frame}

\begin{frame}{Duplication: proof of the theorem}
\framesubtitle{Fact 1 carries many copies up; Fact 2 carries a bad host up}
\small
\keybox{\textbf{Theorem.} $H$ satisfies the conjecture if and only if $R=\operatorname{red}(H)$ does.}
\begin{itemize}
\item<1-> \textbf{If $R$ satisfies it, so does $H$.} Let $\dist_H(A)\ge\epsilon$. $R$ keeps some rows and columns of $H$, so $H$ contains $R$, and every $R$-free matrix is $H$-free: $\dist_R(A)\ge\epsilon$. So $A$ has at least $c\,\epsilon^Cn^{k_0+l_0}$ copies of $R$. Fact 1, once for each repeated row or column, gives at least $c'\epsilon^{C'}n^{k+l}$ copies of $H$. This needs $n$ larger than a fixed power of $1/\epsilon$; for smaller $n$, use $N_H(A)\ge1$ and a larger $C'$.
\item<2-> \textbf{If $R$ fails, so does $H$.} Let $c',C'>0$ be given. As $R$ fails, some $A$ has $\dist_R(A)\ge\epsilon_0$ and fewer than $c'2^{-C'}\epsilon_0^{C'}n^{k_0+l_0}$ copies of $R$. For large $q$, Fact 2 gives
\begin{align*}
\dist_H(A^{(q)})&\ge\tfrac12\dist_R(A)\ge\tfrac{\epsilon_0}2=:\epsilon &&\text{by (b)},\\
N_H(A^{(q)})&\le\frac{N_R(A)}{n^{k_0+l_0}}(qn)^{k+l}<c'\epsilon^{C'}(qn)^{k+l} &&\text{by (a)}.
\end{align*}
So no constants $c',C'$ work for $H$.
\end{itemize}
\end{frame}

\begin{frame}{Infinite families}
\textbf{Counterexamples}
\begin{itemize}
\item $H\oplus(1)\oplus(1)\oplus\cdots$ for OpenAI's $H$: a counterexample of every size from $66$ on.
\item $H$ with zero columns appended: $66\times67$, $66\times68$, \dots
\item If Shea's claim holds: $K=\smallmat{0&1&1\\1&0&1\\1&1&0}$ gives $\smallmat{0&1&1&0\\1&0&1&0\\1&1&0&0\\0&0&0&1}$ and $\smallmat{0&1&1&0\\1&0&1&0\\1&1&0&0}$.
\item Any of these with rows or columns repeated.
\end{itemize}
\medskip
\textbf{Patterns that satisfy the conjecture}
\begin{itemize}
\item Two all-ones blocks in opposite corners, zeros elsewhere: they reduce to $I_2$ (or its mirror image), which satisfies it by Lupini's theorem.
\end{itemize}
\end{frame}

\begin{frame}{Formal verification}
\begin{itemize}
\item \textbf{Lean 4} (Mathlib v4.35.0): the direct-sum theorem, the zero-column variant, their corollaries and the infinite family obtained by repeatedly adding $(1)$ are fully proved: no \texttt{sorry}, only the standard axioms. Written with Kimi 3.1 Pro; checked independently.
\item \textbf{Agda} (\texttt{--safe}): the special case of one appended zero column, fully proved.
\item \textbf{Computer checks}: the direct-sum statements on $200$ small hosts; OpenAI's counts for small $h$, including $m^2$ as the exact distance for $h\le5$.
\end{itemize}
\end{frame}

\begin{frame}{Lean: the definitions}
\framesubtitle{The paper's objects in Lean~4, from \texttt{lean/OrderedRemoval/Statements.lean}}
%%LEAN_DEFS%%

\smallskip
{\scriptsize
\begin{tabular}{@{}l@{\hspace{0.6em}}l@{\hspace{1.4em}}l@{\hspace{0.6em}}l@{}}
$k\times l$ pattern $H$ & \texttt{H : BMat k l}, \texttt{true} $=1$ &
$\smallmat{H&0\\0&Z}$, $\widehat Z$ & \texttt{blockDiag H Z}, \texttt{blowup Z n}\\
$N_H(A)$, $D_H(A)$ & \texttt{copyCount H A}, \texttt{editDist H A} &
no all-zero row or column & \texttt{NoZeroLine H}\\
the conjecture for $H$ & \texttt{PolyRemoval H} &
rows of $Z$ pairwise distinct & \texttt{Function.Injective Z}
\end{tabular}}
\end{frame}

\begin{frame}{Lean: the direct-sum theorem and its corollary}
\framesubtitle{Theorem~8 and Corollary~9 of the paper, as stated in Lean}
%%LEAN_DIRECTSUM%%

\smallskip
\footnotesize
\begin{itemize}
\item \textbf{Also proved:} the copy bound \texttt{copyCount\_blockDiag}, the zero-column variant, and the family from attaching $(1)$ again and again. Step~3's random picks become sums over all picks: no probability or division.
\item \textbf{Lean caught three false edge cases} in the first statements: the empty $0\times0$ pattern satisfies \texttt{NoZeroLine} vacuously, and \texttt{editDist} is \texttt{sInf} $\emptyset=0$ when no $H$-free matrix exists. Fixed by the hypotheses \texttt{hk~:~0~<~k} and ``an $H$-free matrix exists''; the corollaries needed no change.
\end{itemize}
\end{frame}

\begin{frame}{Lean: unconditional, with OpenAI's proof}
\framesubtitle{\texttt{lean-unconditional/}: OpenAI's Lean proof of the $66\times66$ counterexample, plugged into ours}
%%LEAN_UNCOND%%

\medskip
\begin{columns}[T]
\begin{column}{0.54\textwidth}
\footnotesize
\begin{itemize}
\item OpenAI's $%%LEAN_OAI_FILES%%$ files (\texttt{openai/math} at \texttt{%%LEAN_OAI_COMMIT%%}, Apache-2.0), unchanged: we checked all $%%LEAN_OAI_FILES%%$ are byte-identical.
\item A bridge of $%%LEAN_BRIDGE%%$ theorems: both define the same distances and counts, and OpenAI's bound becomes \texttt{$\neg$ PolyRemoval fixedH}.
\item Lean, Mathlib v4.35.0 pinned. An axiom audit (only \texttt{propext}, \texttt{Classical.choice}, \texttt{Quot.sound}) passes on all $%%LEAN_AUDIT_DECLS%%$ declarations; no \texttt{sorry}.
\item Not formalized: the duplication theorem.
\end{itemize}
\end{column}
\begin{column}{0.43\textwidth}
\centering
\begin{tikzpicture}[>={Stealth[length=2mm]}]
\node[leanbox=3.1cm, fill=teal!10] (oai) at (0,0) {OpenAI: $%%LEAN_OAI_FILES%%$ files\\ {\tiny\texttt{no\_polynomial\_removal\_bound}}};
\node[leanbox, fill=accent!12] (ours) at (3.05,0) {ours: $%%LEAN_STMT_LINES%%$ lines\\ {\tiny\texttt{infinite\_family\_of\_}}\\[-1pt] {\tiny\texttt{counterexample}}};
\node[leanbox, fill=paper] (bridge) at (0,-1.35) {bridge: $%%LEAN_BRIDGE%%$ theorems\\ {\tiny\texttt{oai\_not\_polyRemoval}}};
\node[leanbox=4.6cm, fill=green!10] (main) at (1.52,-2.55) {{\tiny\texttt{infinitely\_many\_counterexamples}}};
\draw[->] (oai) -- (bridge);
\draw[->] (bridge) -- (bridge |- main.north);
\draw[->] (ours) -- (ours |- main.north);
\end{tikzpicture}
\end{column}
\end{columns}
\end{frame}

\begin{frame}{Open questions}
\begin{itemize}
\item \textbf{$2\times2$.} The identity $I_2$ satisfies the conjecture (Lupini). The corner $\smallmat{1&0\\1&1}$, OpenAI's body $P$, is open.
\item \textbf{$3\times3$.} Shea's claimed counterexample awaits review.
\item \textbf{Which patterns fail?} Our reductions move counterexamples up in size; what is the smallest one?
\end{itemize}
\end{frame}

\begin{frame}{How we worked}
\small
\begin{itemize}
\item \textbf{Literature, extensions, smaller counterexamples, proofs}: Opus 5.5 (Max), Fable 5.2 (Medium).
\item \textbf{Formalization}: planning with GPT-6.1 Sol (Extra High) and Opus 5.5 (Max); the Lean proof with Open Math Kimi 3.1 Pro (Ultramath); the Agda special case with Opus 5.5 (Max).
\item \textbf{Proofreading, slides, examples}: Opus 5.5 (Max), Open Math DeepSeek V4.1 Flash and Kimi 3.1 Pro.
\end{itemize}
\end{frame}

\begin{frame}{References}
\scriptsize
\begin{itemize}
\item N.~Alon, E.~Fischer, I.~Newman, \emph{Efficient testing of bipartite graphs for forbidden induced subgraphs}, SIAM J. Comput. 37 (2007).
\item N.~Alon, O.~Ben-Eliezer, E.~Fischer, \emph{Testing hereditary properties of ordered graphs and matrices}, FOCS 2017.
\item N.~Alon, O.~Ben-Eliezer, \emph{Efficient removal lemmas for matrices}, Order 37 (2020).
\item L.~Gishboliner, A.~Shapira, \emph{Polynomial property testing}, Comput. Sci. Rev. 58 (2025).
\item OpenAI, \emph{Polynomial removal fails for ordered binary matrices}, preprint, 25 September 2026.
\item ProveIt project, \emph{A $5\times5$ obstruction to polynomial ordered matrix removal}, draft, 7 October 2026.
\item Erebos Labs, \emph{ordered-matrix-removal-4x4}, Lean formalization, 9 October 2026.
\item Z.~Wang, \emph{A four-by-four obstruction to polynomial removal for ordered binary matrices}, arXiv:2610.15370.
\item R.~Shea, \emph{A $3\times3$ ordered binary matrix removal obstruction}, manuscript, 12 October 2026.
\item M.~Lupini, \emph{Polynomial removal for the ordered two-by-two identity matrix}, Zenodo, November 2026.
\end{itemize}
\end{frame}

\end{document}
